\section{结论}
本文围绕 UCI HAR 的跨受试者监督分类，构建人工特征与时序表示的双分支融合模型，并与传统分类器和单分支网络比较。统一外层分组下，逻辑回归取得最高开发均值；PCA 在固定分类器参数的对照中减少维数，但未提高验证 F1。融合的三次基础种子五折均值平均为 0.9314，配置冻结后在官方独立测试上达到 95.55\% 准确率和 0.9564 宏平均 F1，比逻辑回归净减少 17 个错误。

融合收益主要体现在动态活动，坐与站立的 F1 反而下降，部分受试者也未受益。因此，模型复杂度不能替代与任务一致的评价设计。本文以分组验证、独立测试和可追溯结果留档支持有限的实证结论；后续应在新评价数据和更对称的模型选择协议下验证其可推广性。
